\clearpage
\chapter*{ABSTRAK}
\addcontentsline{toc}{chapter}{ABSTRAK}
\begin{center}
    \center
    \begin{singlespace}
        \large\bfseries\MakeUppercase{\thetitle}

        \normalfont\normalsize
        Oleh:

        \bfseries \theauthor
    \end{singlespace}
\end{center}

\begin{singlespace}
\small
Sistem Rekam Medis Elektronik (RME) konvensional menyimpan data kesehatan secara tersebar pada setiap fasilitas pelayanan kesehatan sehingga tenaga kesehatan sulit memperoleh gambaran kondisi pasien secara utuh. 
DecMed (Decentralized Medical Record Management System) telah dikembangkan dengan pendekatan \textit{patient-centric} untuk mengintegrasikan data tersebut dan memberikan kendali data kepada pasien. 
Namun, sumber data DecMed masih bersifat \textit{provider-centric} karena hanya mengakomodasi data dari tenaga kesehatan melalui kunjungan klinis, belum mencakup \textit{Patient-Generated Health Data} (PGHD) yang dihasilkan pasien secara mandiri melalui perangkat \textit{wearable}, sensor kesehatan, \textit{smartphone}, maupun pencatatan manual. 
Data tersebut dapat melengkapi gambaran kondisi pasien di antara kunjungan klinis dan mendukung pengambilan keputusan klinis, sehingga diperlukan modul integrasi PGHD pada DecMed.

Penelitian ini mengembangkan modul integrasi PGHD pada DecMed yang mengumpulkan dan mengirimkan data \emph{time-series} secara \emph{offline-first}, mempertahankan data saat perangkat luring, melakukan sinkronisasi setelah koneksi tersedia, mencatat \emph{data provenance}, serta menjaga kerahasiaan dan integritas data melalui \emph{client} Android pasien, \emph{Proxy Re-Encryption} Server, \emph{smart contract} IOTA, dan IPFS\@.

Pengujian menunjukkan sistem berhasil memproses 151.200 \textit{record} PGHD sintetis pada laju 600 \textit{record} per detik tanpa kehilangan data maupun \textit{crash}, serta menerima 20 \textit{batch} konkuren dari 20 pasien tanpa kegagalan. Variasi manipulasi terhadap \textit{ciphertext}, hash, tanda tangan digital, referensi CID, dan metadata berhasil dideteksi dan ditolak, \textit{plaintext} tidak ditemukan pada titik yang diperiksa, dan akses tanpa persetujuan sah berhasil ditolak. Data tetap tersimpan saat perangkat luring, pengiriman dijadwalkan ulang, dan diproses kembali setelah koneksi tersedia tanpa kehilangan data, sehingga PGHD terbukti dapat diintegrasikan ke DecMed melalui modul ini sesuai skenario pengujian.

\textbf{\textit{Kata kunci: \textit{Data Provenance}, DecMed, Integritas, Kerahasiaan, Ketersediaan, \textit{Patient-Generated Health Data}, Rekam Medis Elektronik}}
\end{singlespace}
\clearpage